\chapter{DevOps, GitOps, dan Platform Engineering}

\section*{Tujuan Pembelajaran}
\begin{itemize}
  \item Menjelaskan praktik DevOps dan metrik DORA (DevOps Research and Assessment) untuk mengukur kinerja pengiriman perangkat lunak.
  \item Menjelaskan alur CI/CD (\textit{Continuous Integration / Continuous Delivery}) dan empat prinsip GitOps.
  \item Merancang topologi tim dan peran platform internal bagi organisasi.
\end{itemize}

\section{Pendahuluan}

% TODO
% Cakupan: satu perubahan kecil di Medikara Health butuh tiga bulan karena harus melewati banyak tiket dan persetujuan manual; arsitektur cloud-native tidak memberi manfaat bila cara mengirim perubahan tetap lambat.

\section{DevOps dan Metrik DORA}

% TODO
% Cakupan: DevOps sebagai budaya dan praktik; metrik DORA (frekuensi deployment, lead time perubahan, tingkat kegagalan perubahan, waktu pemulihan); temuan laporan DORA 2025 bahwa AI memperkuat kekuatan maupun kelemahan organisasi; anggapan keliru bahwa rilis yang jarang lebih aman.
% Kasus Medikara: perkiraan metrik DORA kondisi saat ini untuk SIMRS dan Medikara+.

\section{CI/CD dan Strategi Rilis}

% TODO
% Cakupan: continuous integration, continuous delivery, dan continuous deployment; tahapan pipeline; strategi rilis blue-green dan canary; keamanan rantai pasok perangkat lunak disinggung singkat (dibahas rinci di IT40813).
% Pendalaman teknis: membuat pipeline CI sederhana dengan GitHub Actions.

\section{Infrastructure as Code dan GitOps}

% TODO
% Cakupan: IaC sebagai infrastruktur yang ditulis dan ditinjau seperti kode; empat prinsip OpenGitOps v1.0.0 (declarative, versioned and immutable, pulled automatically, continuously reconciled); jejak audit yang dihasilkan GitOps.
% Pendalaman teknis: menyinkronkan manifest Kubernetes dengan Argo CD atau Flux.

\section{Platform Engineering dan Internal Developer Platform}

% TODO
% Cakupan: platform engineering dan IDP menurut CNCF Platforms White Paper; golden path, self-service, dan platform as a product; Platform Engineering Maturity Model (investasi, adopsi, antarmuka, operasi, pengukuran); perbedaan tim platform dengan divisi infrastruktur yang melayani tiket.

\section{Team Topologies}

% TODO
% Cakupan: empat tipe tim (stream-aligned, platform, enabling, complicated-subsystem) dan tiga mode interaksi (collaboration, X-as-a-Service, facilitating) menurut Skelton dan Pais edisi ke-2 (2025); beban kognitif tim; kaitan dengan hukum Conway (Bab 7).
% Kasus Medikara: penyatuan tim TI dari lima rumah sakit menjadi tim stream-aligned dan tim platform.

\section{Aktivitas: Model Operasi Delivery dan Peta Jalan Platform Internal}

% TODO
% Tujuan aktivitas: memetakan alur perubahan saat ini, merancang alur CI/CD dan GitOps target, serta menyusun topologi tim dan kapabilitas platform internal yang diprioritaskan.
% Peran dalam tim: anggota non-teknis memimpin pemetaan waktu tunggu dan rantai persetujuan; anggota teknis merancang alur CI/CD dan GitOps target; seluruh anggota menyusun topologi tim dan prioritas platform.
% Langkah: (1) petakan alur perubahan dari permintaan fitur sampai produksi; (2) tandai antrean dan langkah manual; (3) rancang alur target; (4) susun topologi tim target; (5) pilih tiga sampai lima kapabilitas platform internal dan urutannya.
% Luaran dan tonggak capstone: C4, model operasi delivery dan peta jalan platform internal.
% Berkas sources yang direncanakan: sources/chapter13/templat-pemetaan-alur-perubahan.xlsx, sources/chapter13/templat-topologi-tim.drawio, sources/chapter13/contoh-pipeline-gitops/ (opsional)

\section{Ringkasan}

% TODO

\section{Latihan}

\begin{enumerate}
  \item Organisasi yang merilis perubahan setiap tiga bulan sering menganggap rilis yang jarang lebih aman. Kritisi anggapan ini dengan merujuk temuan DORA.
  \item Apa perbedaan tim platform dengan divisi infrastruktur tradisional yang melayani tiket?
  \item Jelaskan cara GitOps membantu auditor menelusuri siapa mengubah apa di lingkungan produksi.
\end{enumerate}
